\documentclass[conference]{IEEEtran}
\IEEEoverridecommandlockouts
% The preceding line is only needed to identify funding in the first footnote. If that is unneeded, please comment it out.
\usepackage{cite}
\usepackage{amsmath,amssymb,amsfonts}
\usepackage{algorithmic}
\usepackage{graphicx}
\usepackage{textcomp}
\usepackage{xcolor}
\def\BibTeX{{\rm B\kern-.05em{\sc i\kern-.025em b}\kern-.08em
    T\kern-.1667em\lower.7ex\hbox{E}\kern-.125emX}}
\begin{document}

\title{AANEC-Relaxed Trapped Surface Theorem and Stellar Mass-Gap Derivation}\author{Anonymous Research Plan Author}\maketitle

\begin{abstract}
This proposal develops a qualified AANEC-relaxed trapped-surface theorem for classical general relativity under spherical, isotropic stellar collapse. The central argument is that pointwise energy conditions are too restrictive for realistic quantum and semiclassical matter, so the formal route replaces them with the Averaged Null Energy Condition plus an explicit, independent averaged null-convergence/focusing condition. The planned contribution is not a proof but a structured conjecture: if global hyperbolicity, regular matter, complete null generators, and this added focusing hypothesis hold, then exceeding a model-dependent Misner-Sharp compactness threshold should force trapped-surface formation, thereby excluding stable non-black-hole remnants above a compactness-dependent mass scale within the restricted class. The observed 2-5 solar-mass gap is treated as an astrophysical consistency constraint, not as a theorem-derived prediction, because explosion energetics, fallback, equation-of-state physics, rotation, magnetic fields, nonspherical dynamics, and binary evolution remain viable alternative explanations.

| Decision | Condition | Consequence | |---|---|---| | Retain theorem | AANEC plus explicit averaged focusing holds on complete null generators | Trapped-surface threshold is proposed and remnant exclusion follows only in restricted model class | | Restrict or reject theorem | Averaged focusing fails, generators incomplete, quantum ANEC violations, or non-spherical dynamics dominate | The compactness bound is not derivable from stated assumptions | | Demote mass-gap inference | Astrophysical mechanisms reproduce remnant distribution without compactness bound | Gap remains consistent with theorem but loses discriminating power for this mechanism |
\end{abstract}
\begin{IEEEkeywords}
AANEC, trapped surface, compactness threshold, stellar mass gap, global hyperbolicity, null generators
\end{IEEEkeywords}\section{Introduction}
\subsection{From Pointwise Energy Conditions to an Averaged Focusing Hypothesis}
Classical trapped-surface arguments traditionally rest on pointwise energy conditions, yet quantum fields and non-minimally coupled matter systematically violate such pointwise positivity. The literature on energy conditions in general relativity and quantum field theory motivates a shift toward averaged statements, notably the achronal averaged null energy condition (AANEC), as a more robust geometric hypothesis for singularity theorems. This proposal advances a route-specific contribution by treating AANEC not as a sufficient condition, but as one ingredient in a qualified mechanism: AANEC plus an explicit averaged null-convergence/focusing condition along complete null generators. The latter is declared as an independent geometric hypothesis, not a consequence of AANEC alone, thereby separating the formal mechanism from the broader claim that averaged energy conditions universally imply focusing. The resulting theorem/conjecture is scoped to classical general relativity with spherical, nonrotating collapse and isotropic pressure; it is not a quantitative derivation of the observed 2-5 solar mass gap, nor a proof that all massive stars must form black holes. See Eq.~\eqref{eq:introduction-introduction-002}.~\cite{cite_0ed84243cf171bc4}

\subsection{AANEC as a Necessary but Not Sufficient Focusing Ingredient}
\begin{equation}
int_{gamma} T_{ab} k^{a} k^{b} dlambda >= 0
\label{eq:introduction-introduction-002}
\end{equation}

\subsection{Separating the Formal Bound from the Astrophysical Mass-Gap Inference}
The compactness threshold that forces trapped-surface formation is model-dependent and requires a precise definition of the Misner-Sharp parameter; until that symbol-domain convention is fixed, the threshold is stated only as a proposed dependency rather than a computed number. The proposed contribution therefore keeps three claims distinct. First, the formal conjecture links averaged focusing and compactness to trapped-surface existence within the restricted model class. Second, the exclusion of stable non-black-hole remnants above a compactness-dependent scale depends on a TOV-like stability predicate that must itself be formalized. Third, the observed lower stellar mass gap is treated as an astrophysical consistency constraint, not as theorem-derived evidence, because explosion energetics, fallback, equation-of-state physics, pair-instability processes, rotation, magnetic fields, nonspherical dynamics, and binary evolution can populate or suppress the gap independently of any trapped-surface bound. Gravitational-wave population analyses and binary-evolution studies sharpen this separation: detections and inferred distributions constrain formation channels without establishing the formal theorem, while events that challenge current compact-object models indicate where the mass-gap inference becomes non-discriminating. The next stage must therefore fix the compactness and stability conventions before the Raychaudhuri/integrated-threshold chain can be developed into a verifiable proposition. Specifically, the definitions ledger must resolve the Misner-Sharp parameter and the averaged focusing condition to enable the forward derivation. See Eq.~\eqref{eq:introduction-introduction-002}.~\cite{cite_72d7ff391b29c495,cite_c369df652181e8f7,cite_7a407c865faafd11,cite_fd0f77b2a173c0ab,cite_977568cf7b8a55bf,cite_07d6df5d6978a1fc,cite_f634820e22d7a238,cite_cbbced417410bc6c,cite_42e341ff6471ce89,cite_7a44e6730c44cb39}

\begin{table*}[htbp]
\centering
\footnotesize
\setlength{\tabcolsep}{3pt}
\renewcommand{\arraystretch}{1.12}
\caption{Route-Specific Contribution and Scope Decisions}
\label{tab:introduction-introduction-004}
\begin{tabular}{|p{0.29\textwidth}|p{0.29\textwidth}|p{0.29\textwidth}|}
\hline
Decision point & Adopted scope for this route & Consequence if ignored \\\hline
Energy-condition premise & AANEC plus an explicit averaged null-convergence/focusing hypothesis & AANEC alone may not imply the integrated focusing needed for trapped-surface formation \\\hline
Compactness threshold & Model-dependent; formal Misner-Sharp definition deferred to the definitions ledger & Thresholds and counterexample tests remain underdetermined \\\hline
Mass-gap status & Astrophysical consistency constraint, not a theorem prediction & Explosion, fallback, equation-of-state, rotation, magnetic, nonspherical, or binary mechanisms can mimic or erase the gap \\\hline
Counterexample criterion & Witness must satisfy all declared assumptions, including the added focusing condition & Witnesses that evade the bound by violating the extra hypothesis do not falsify the conjecture \\\hline
\end{tabular}
\end{table*}~\cite{cite_72d7ff391b29c495,cite_c369df652181e8f7,cite_7a407c865faafd11,cite_fd0f77b2a173c0ab,cite_977568cf7b8a55bf,cite_07d6df5d6978a1fc,cite_f634820e22d7a238,cite_cbbced417410bc6c,cite_42e341ff6471ce89,cite_7a44e6730c44cb39}

\section{Background, Survey, and Research Gap}
\subsection{Theoretical Foundations and the Energy Condition Bridge}
The formal argument proposed in this study is motivated by a specific tension within classical general relativity: the reliance of Penrose-type singularity theorems on pointwise energy conditions, which are systematically violated by realistic quantum fields. The survey evidence establishes that pointwise positivity is incompatible with quantum mechanics, as phenomena such as the Casimir effect and squeezed vacuum states demonstrate local negative energy densities that violate the Null Energy Condition (NEC) and Strong Energy Condition (SEC) [W3009127927]. This violation challenges the robustness of singularity theorems, which traditionally require the NEC to ensure the focusing of null geodesics via the Raychaudhuri equation. To preserve the predictive power of these theorems in semiclassical regimes, the theoretical landscape has shifted toward averaged conditions. The Achronal Averaged Null Energy Condition (AANEC) has emerged as a candidate for a universal property of gravitating matter, offering a pathway to maintain geodesic incompleteness predictions despite local quantum violations [W3009127927]. However, a critical unresolved bridge remains: the survey indicates that while AANEC is a weaker assumption than pointwise NEC, it is not established that AANEC alone implies the specific local or integrated focusing required for trapped-surface formation in dynamic stellar collapse. This gap motivates the proposed addition of an explicit averaged null-convergence condition as an independent geometric hypothesis, the formalization of which is owned by the definitions ledger.~\cite{cite_0ed84243cf171bc4}

\subsection{Astrophysical Context: The Stellar Mass Gap as a Consistency Constraint}
Parallel to the theoretical foundations, the survey synthesizes observational evidence regarding the lower stellar mass gap, a region of relative scarcity of compact objects between approximately 2 and 5 solar masses. This gap serves as a stringent discriminator for supernova engine physics and collapse outcomes. Dynamical mass measurements in X-ray binaries and gravitational-wave population analyses confirm a robust absence of objects in this interval, suggesting a discontinuity in the formation of neutron stars versus black holes [W2994776863][W3036761016]. The survey attributes this feature primarily to astrophysical mechanisms rather than pure general-relativistic bounds: specifically, the interplay between progenitor compactness, explosion energetics, and fallback dynamics. For instance, models indicate that rapid instabilities in core collapse can lead to either successful explosions (leaving neutron stars) or direct collapse (leaving black holes), thereby avoiding the intermediate mass range [W2994776863]. Furthermore, the location of the lower edge of the gap is dictated by near-final core mass and is robust against variations in metallicity and wind mass loss [W2994776863]. However, the presence of anomalous detections, such as the 2.6 solar mass object in GW190814, challenges simple models and highlights the complexity of the boundary between neutron-star and black-hole formation [W3036761016]. These observations define the astrophysical consistency constraints that the proposed formal theorem must address without claiming to replace stellar-evolution explanations.~\cite{cite_7a407c865faafd11}

\subsection{Research Gap: The Missing Formal Claim for AANEC-Relaxed Trapped Surfaces}
The primary research gap identified by the survey is the absence of a rigorous formal claim establishing the validity domain of Penrose-type singularity theorems when applied to AANEC-relaxed assumptions in the context of stellar collapse. Specifically, the survey evidence plan flags a missing 'formal\_claim' slot for the sub-hypothesis concerning whether AANEC plus an explicit averaged focusing condition yields a trapped-surface bound [survey:evidence\_plan]. Current literature discusses AANEC as a necessary condition for singularity theorems in quantum regimes, but does not provide a complete derivation of a compactness threshold for trapped-surface formation under the specific combination of AANEC and an added averaged null-convergence hypothesis [W3009127927]. This gap is critical because without such a formal claim, the connection between the geometric relaxation of energy conditions and the astrophysical mass gap remains speculative. The proposed study aims to fill this gap by deriving a qualified trapped-surface theorem that explicitly incorporates an averaged focusing condition distinct from standard AANEC, thereby clarifying the assumptions required to exclude stable non-black-hole remnants above a certain compactness.

\subsection{Research Gap: Unmapped Boundary Variables in Stellar Evolution}
A secondary research gap involves the identification of precise boundary variables that separate black-hole from non-black-hole outcomes in theoretical stellar evolution models. The survey notes that the 'boundary\_variable' slot remains uncovered in the relevant evidence cluster 'Theoretical Stellar Evolution and Remnant Formation Pathways' [survey:evidence\_plan]. While post-bounce compactness metrics (e.g., xi\_2.5) and explosion timescales are recognized as discriminators, the survey highlights that the specific formal variable governing the transition under AANEC-relaxed conditions is not defined in the existing literature [W2808391388]. This lack of a mapped boundary variable prevents a direct comparison between the proposed geometric compactness threshold and the astrophysical parameters that determine remnant mass. The gap is further complicated by the sensitivity of mass distributions to uncertainties in explosion and unbinding energies, which remain unconstrained across different progenitor masses [survey:gap\_ledger]. Consequently, the research plan must treat the astrophysical boundary variable as a distinct, unresolved parameter that interacts with, but is not derived from, the formal geometric bound.

\subsection{Alternative Explanations and Confounders}
The survey emphasizes that the observed mass gap and compactness thresholds are subject to significant alternative explanations that must be separated from the formal theorem. These confounders include rotation, magnetic fields, nonspherical dynamics, binary evolution history, and equation-of-state uncertainties [W2149573451][W2980421388]. For example, numerical simulations show that the use of non-isentropic equations of state can lead to delayed collapse to a black hole due to pressure support from shock heating, altering the compactness threshold compared to isentropic models [W2149573451]. Similarly, binary evolution channels, including common envelope phases and mass transfer, can populate or avoid the mass gap independently of the core-collapse physics [W2980421388]. These mechanisms suggest that the observed distribution of remnant masses is not solely a consequence of a general-relativistic trapped-surface bound. Therefore, the research gap includes the need to explicitly demote the observed mass gap from a theorem-derived prediction to an astrophysical consistency constraint, acknowledging that explosion energetics and fallback processes may dominate the observed demographics [survey:gap\_ledger].~\cite{cite_7a44e6730c44cb39}

\begin{table*}[htbp]
\centering
\footnotesize
\setlength{\tabcolsep}{3pt}
\renewcommand{\arraystretch}{1.12}
\caption{Decision Matrix: Evidence Status and Gap Implications}
\label{tab:survey-and-research-gap-gap-04}
\begin{tabular}{|p{0.29\textwidth}|p{0.29\textwidth}|p{0.29\textwidth}|}
\hline
Evidence Category & Survey Status & Implication for Proposed Theorem \\\hline
**Energy Conditions** & AANEC is a candidate universal property; pointwise NEC violated by quantum fields. & Theorem must relax pointwise NEC to AANEC but requires an added explicit averaged focusing condition to guarantee trapped surfaces. \\\hline
**Mass Gap Observations** & Robust scarcity in 2-5 solar mass range; anomalies like GW190814 challenge simple models. & Theorem cannot predict the exact mass gap; it must provide a compactness bound that is consistent with, but not identical to, the observed gap. \\\hline
**Stellar Evolution** & Boundary variables (compactness, explosion energy) are identified but not formally mapped to GR bounds. & The formal compactness threshold is a theoretical construct; its physical interpretation requires separate astrophysical inference. \\\hline
**Counterexamples** & Regular black holes and exotic compact objects evade classical singularity theorems. & The theorem's validity domain must explicitly exclude or address these alternatives via the averaged focusing assumption. \\\hline
\end{tabular}
\end{table*}

\section{Research Questions and Planned Contributions}
\subsection{Scope and Research Questions}
This route investigates whether replacing pointwise energy conditions with the Averaged Null Energy Condition (AANEC) and an explicit averaged null-convergence assumption yields a rigorous trapped-surface bound for spherical stellar collapse. The central research question asks: under what precise geometric and matter assumptions does a model-dependent compactness threshold force trapped-surface formation, and what does this imply for stable non-black-hole remnants? A secondary question addresses the scope of this formal result relative to astrophysical observations, specifically whether the derived bound explains the observed 2-5 solar mass gap or merely provides a consistency constraint against stellar-physics confounders such as explosion energetics, fallback, and binary evolution.~\cite{cite_0ed84243cf171bc4,cite_7a407c865faafd11,cite_f634820e22d7a238}

\subsection{Planned Contributions and Theoretical Mechanism}
The primary contribution is a proposed theorem linking AANEC satisfaction to trapped-surface formation via Raychaudhuri-type focusing, contingent on a distinct averaged null-convergence condition. This mechanism replaces standard pointwise assumptions, which are often violated in quantum regimes, with integrated constraints along null generators. The plan involves deriving a compactness threshold where outgoing and ingoing null expansions become negative, thereby excluding stable non-black-hole remnants above a critical mass scale within the restricted model class. This formal contribution is designed to operate independently of specific supernova engine models, treating the mass gap as an astrophysical consistency check rather than a direct theorem prediction.

\subsection{Proposed Averaged Null-Convergence Condition}
\begin{equation}
The formalization of the explicit averaged null-convergence condition (A5) is a critical dependency. While the exact mathematical form distinguishing it from standard AANEC (A2) requires human input, the proposed structure involves an integral constraint on the stress-energy tensor along null geodesics k^a. The condition is hypothesized to ensure sufficient focusing of the congruence to trigger the Raychaudhuri equation's convergence term, distinct from the pointwise positivity of the Null Energy Condition.
\label{eq:research-questions-and-contributions-block-3}
\end{equation}

\subsection{Unresolved Items and Delimiting Questions}
The operationalization of these contributions is delimited by several unresolved formal definitions owned by the definitions ledger. These gaps determine the boundaries of the proposed theorem and the validity of candidate counterexamples. Consequently, the research questions must be interpreted as conditional inquiries: they ask what follows if and only if the specific geometric and stability predicates are formally resolved. This section advances the proposal by framing the inquiry around these dependencies, ensuring that the planned contributions target the logical gaps identified in the survey rather than assuming their resolution. See Eq.~\eqref{eq:research-questions-and-contributions-block-3}.

\section{Idea Source Checkpoints and Direction Selection Audit}
\subsection{Scientific Attraction and Mechanism Replacement}
The retained direction, mechanism replacement, advances a formal conjecture for classical general relativity: if spherically symmetric stellar collapse satisfies the Averaged Null Energy Condition (AANEC), global hyperbolicity, regular matter fields, complete null generators, and an explicit averaged null-convergence/focusing condition, then reaching a model-dependent compactness threshold forces trapped-surface formation. Consequently, stable non-black-hole remnants above a compactness-dependent mass scale are excluded within this restricted model class. This proposal replaces pointwise energy conditions with AANEC plus an explicit averaged focusing assumption, deriving a trapped-surface threshold via Raychaudhuri-type focusing and Misner-Sharp compactness. The scientific attraction lies in relaxing restrictive pointwise assumptions while maintaining a rigorous geometric bound. However, the observed 2-5 solar mass gap is explicitly demoted to an astrophysical consistency constraint rather than a theorem-derived prediction, acknowledging that supernova explosion physics, fallback, equation-of-state uncertainties, rotation, magnetic fields, nonspherical dynamics, or binary interactions may govern the observed distribution.~\cite{cite_0ed84243cf171bc4}

\subsection{Defects and Unsupported Links}
Upstream checkpoints expose critical defects in the formal reasoning chain, primarily stemming from undefined mathematical constructs owned by the definitions ledger. These missing definitions render the proof obligations unresolved and prevent the closure of the derivation steps. Without these formal definitions, the logical chain from AANEC to trapped-surface formation remains unverified, and the exclusion of stable non-black-hole remnants cannot be rigorously established. This section advances the proposal by identifying these specific logical breaks, thereby justifying the need for a qualified, dependency-aware formalization in the subsequent stages.

\begin{table*}[htbp]
\centering
\footnotesize
\setlength{\tabcolsep}{3pt}
\renewcommand{\arraystretch}{1.12}
\caption{Counterexample Evaluation Matrix}
\label{tab:idea-origin-and-selection-block-3}
\begin{tabular}{|p{0.23\textwidth}|p{0.23\textwidth}|p{0.23\textwidth}|p{0.23\textwidth}|}
\hline
Candidate Witness & Assumption A5 (Averaged Focusing) & Assumption A6 (Stable Non-BH Remnant) & Validity Status \\\hline
CE1: Hypothetical static regular metric with high compactness & False (Likely violates additional A5 constraint to evade trapped surfaces) & False (Definition conflict: stable objects must be below C\_max) & Boundary Case \\\hline
CE2: Regular black hole model (e.g., Hayward/Bardeen) & Unknown (Modifies stress-energy/causal structure) & False (Contains horizons/trapped surfaces, thus classified as BH) & Assumptions Not Satisfied \\\hline
\end{tabular}
\end{table*}

\subsection{Qualified Retained Direction and Transition}
The retained direction is qualified by the explicit exclusion of astrophysical mechanisms as direct proof obligations. The theorem fails or must be restricted if AANEC is violated, the added averaged focusing condition fails, null generators are incomplete, quantum stress-energy violates ANEC, nonspherical/rotating/magnetic dynamics dominate, or stable exotic compact objects evade the averaged-focusing assumptions. The mass-gap inference fails or becomes non-discriminating if explosion energetics, fallback, equation-of-state physics, pair-instability processes, or binary evolution dominate the observed remnant distribution. This section establishes the scope of the formal conjecture and identifies the missing mathematical definitions as the primary barriers to proof closure. The next stage, Definitions and Propositions, must provide the rigorous formalization of these dependencies to enable the forward derivation of the trapped-surface threshold and the evaluation of counterexample validity.~\cite{cite_7a407c865faafd11,cite_f634820e22d7a238}

\section{Problem Definition, Assumptions, and Hypotheses}
\subsection{Formal Problem and Scope}
The central problem is to determine whether spherically symmetric stellar collapse in classical General Relativity, governed by the Averaged Null Energy Condition (AANEC) and an explicit averaged null-convergence/focusing assumption, forces trapped-surface formation above a model-dependent compactness threshold. If such a threshold exists, stable non-black-hole remnants exceeding it would be excluded within the restricted model class. This formal inquiry is distinct from astrophysical explanations for the observed 2-5 solar mass gap, which involve complex stellar physics not captured by the theorem. The scope is strictly limited to regular matter fields, global hyperbolicity, and complete null generators, excluding nonspherical, rotating, or magnetic dynamics.

\subsection{AANEC and Focusing Relations}
int\_\{-∞\}\textasciicircum{}\{∞\} T\_\{ab\} k\textasciicircum{}\{a\} k\textasciicircum{}\{b\} dλ ≥ 0

∫\_\{γ\} R\_\{ab\} k\textasciicircum{}\{a\} k\textasciicircum{}\{b\} dλ > 0~\cite{cite_0ed84243cf171bc4}

\subsection{Interpretation of Energy and Focusing}
The first relation defines the AANEC along a complete null geodesic γ with tangent vector k\textasciicircum{}a and affine parameter λ. The second relation represents the proposed explicit averaged null-convergence/focusing condition, which is treated as an independent geometric hypothesis rather than a direct consequence of AANEC alone. This distinction is critical: while AANEC ensures non-negative averaged stress-energy, the explicit focusing condition is required to drive the expansion scalar negative via the Raychaudhuri equation, thereby enabling trapped-surface formation. The mathematical form of this second condition remains unresolved and is owned by the definitions ledger, serving as a key dependency for the subsequent derivation.~\cite{cite_0ed84243cf171bc4}

\subsection{Assumption Ledger}
\begin{itemize}
\item A1: Classical GR with spherical symmetry and isotropic pressure.
\item A2: Matter fields satisfy AANEC along complete null generators.
\item A3: Spacetime is globally hyperbolic and regular.
\item A4: Null generators are complete.
\item A5: Explicit averaged null-convergence/focusing condition holds (form undefined).
\item A6: Non-black-hole remnants are stable equilibrium objects with bounded compactness.
\end{itemize}

\subsection{Proposed Exclusion Proposition}
\paragraph{Proposition (proposed).} Proposition P2: Stable non-black-hole remnants above a compactness-dependent mass scale are excluded within the restricted model class defined by AANEC and averaged focusing. Proof Obligation PO1 requires establishing that A5 implies negative expansion. Failure Condition: The proposition is invalidated if a counterexample exists satisfying A1-A6 but avoiding trapped surfaces, or if A5 is proven redundant to A2 in this context.

\section{Study Design and Methods}
\subsection{Protocol Scope and Unit of Analysis}
This study adopts a formal proof-and-counterexample search as its unit of analysis, targeting a restricted model class in classical general relativity. The protocol replaces pointwise energy conditions with the Averaged Null Energy Condition (AANEC) and an explicit, independent averaged null-convergence assumption on null congruence generators. This relaxation is motivated by the systematic violation of pointwise conditions in quantum field theory, which necessitates weaker, averaged statements for broader validity [W3009127927]. The design explicitly separates formal geometric consequences from astrophysical inference: the observed stellar mass gap is treated solely as a consistency constraint subject to confounders such as explosion energetics, fallback, and binary evolution, rather than as a direct prediction of the theorem. The protocol does not execute numerical simulations or empirical tests; it defines the logical boundaries within which a trapped-surface theorem and a subsequent stable-remnant exclusion bound may be rigorously evaluated or falsified.~\cite{cite_0ed84243cf171bc4}

\begin{table*}[htbp]
\centering
\footnotesize
\setlength{\tabcolsep}{3pt}
\renewcommand{\arraystretch}{1.12}
\caption{Formal Evaluation and Boundary Decision Matrix}
\label{tab:study-design-and-methods-decision-matrix}
\begin{tabular}{|p{0.23\textwidth}|p{0.23\textwidth}|p{0.23\textwidth}|p{0.23\textwidth}|}
\hline
Evaluation Dimension & Decision Criterion & Required Upstream Artifact & Protocol Consequence \\\hline
Trapped-Surface Formation & Both null expansions negative at threshold & Raychaudhuri focusing lemma & Proposition accepted if proof obligations closed \\\hline
Stable Remnant Exclusion & No equilibrium solution above compactness bound & TOV-like stability predicate & Exclusion bound derived for restricted model class \\\hline
Counterexample Validity & Witness satisfies all independent assumptions & Explicit averaged focusing definition & Candidate rejected if any assumption fails \\\hline
Mass-Gap Relevance & Astrophysical confounders ruled out & Explosion/fallback sensitivity analysis & Gap treated as non-discriminating for theorem \\\hline
\end{tabular}
\end{table*}

\subsection{Comparators, Ablations, and Robustness Checks}
The protocol specifies three ablation comparisons to isolate the contribution of the averaged focusing assumption from standard AANEC. First, the baseline comparison evaluates whether AANEC alone suffices for trapped-surface formation; if it does, the independent focusing assumption is redundant and the theorem reduces to existing Penrose-type results. Second, the counterfactual ablation removes the explicit averaged focusing condition while retaining AANEC, global hyperbolicity, and complete null generators; a valid witness in this configuration would demonstrate that AANEC is insufficient for the claimed bound. Third, the boundary-robustness check tests whether the compactness threshold shifts under controlled relaxations of spherical symmetry or isotropic pressure. These comparisons are formal rather than empirical: they assess logical dependence within the restricted model class. The astrophysical comparators—explosion energy, fallback, and equation-of-state variation—are retained only as confounder checks, not as evidence for or against the geometric theorem.

\subsection{Counterexample and Boundary Analysis Criteria}
A valid counterexample must satisfy all declared assumptions simultaneously: spherical symmetry, isotropic pressure, AANEC, global hyperbolicity, complete null generators, and the explicit averaged focusing condition, while failing to form a trapped surface or admitting a stable non-black-hole remnant above the compactness bound. Candidate witnesses drawn from regular black hole models or ultra-compact object constructions are screened against each assumption in sequence. A witness that evades trapped-surface formation by violating the averaged focusing condition is classified as a boundary case rather than a counterexample, since it falls outside the theorem's scope. A witness that satisfies all assumptions yet evades the bound would falsify the proposition and necessitate restriction of the model class. The screening procedure is purely formal; it does not adjudicate the physical plausibility of the witness metric, only its logical consistency with the stated premises.

\subsection{Artifacts, Reproducibility, and Required Human Decisions}
The formal artifacts—assumption ledger, proof obligations, and counterexample screening records—require a data governance and reproducibility protocol that has not yet been specified. The current plan lacks a finalized data management strategy for versioning formal reasoning steps and a reproducibility standard for independent verification of the proof chain. Additionally, the measurement and calibration plan for the astrophysical consistency checks remains undefined: no measurement protocol or quality-control procedure has been established for the remnant mass distributions used as confounder controls. These gaps are flagged as required human decisions before the formal search can proceed to verification. The absence of these governance and measurement specifications does not invalidate the formal protocol but prevents its execution under standard reproducibility criteria. This section advances the proposal by explicitly bounding the methodological scope, ensuring that the formal argument is evaluated independently of unresolved empirical calibration issues.

\section{Expected Outcome Branches and Conditional Conclusions}
\subsection{Interpretive Scope of the Formal Design}
This section maps the formal proof-and-counterexample search to its predeclared outcome branches. The experimental unit is not an observational survey but a decision procedure over the scoped propositions: a proposed trapped-surface threshold under averaged focusing, and the consequent exclusion of stable non-black-hole remnants above a compactness-dependent mass scale. Because the design is mathematical, every branch below is an expected result of the analysis, not an observed finding. The central interpretive rule is that the theorem's validity domain is the product of its assumptions, so a branch outcome must be read as a statement about the restricted model class rather than about all astrophysical collapse. This approach ensures that the outcomes stage consumes the definitions ledger's formalizations without repeating their detailed status.

\subsection{Decision Criterion Linking Focusing and Compactness}
F\_\{\textbackslash{}lambda\}(k) \textbackslash{}equiv \textbackslash{}int\_\{\textbackslash{}gamma\_\{k\}\} R\_\{ab\}k\textasciicircum{}\{a\}k\textasciicircum{}\{b\}\textbackslash{},d\textbackslash{}lambda \textbackslash{}geq 0 \textbackslash{}quad \textbackslash{}wedge \textbackslash{}quad C \textbackslash{}geq C\_\{\textbackslash{}mathrm\{max\}\} \textbackslash{};\textbackslash{}Longrightarrow\textbackslash{}; \textbackslash{}exists\textbackslash{},\textbackslash{}text\{trapped surface\} \textbackslash{};\textbackslash{}Longrightarrow\textbackslash{}; \textbackslash{}neg\textbackslash{},\textbackslash{}text\{stable non-BH remnant\}

\subsection{Reading the Criterion as a Branch Selector}
The displayed relation is the decision spine that each branch tests. It conjoins the proposed averaged null-convergence condition with the compactness threshold to force trapped-surface formation, which in turn contradicts the stable-remnant predicate. The left-hand side is the independent geometric hypothesis that the design adds to the averaged energy condition; the right-hand side is the exclusion claim under audit. A branch outcome is therefore determined by whether the analysis closes the obligations attaching to the focusing term, the threshold term, and the stability predicate, or instead exposes a failure in one of them. This criterion is proposed and unverified; it is not a completed derivation.

\begin{table*}[htbp]
\centering
\footnotesize
\setlength{\tabcolsep}{3pt}
\renewcommand{\arraystretch}{1.12}
\caption{Conditional-Outcome Decision Matrix}
\label{tab:expected-outcomes-eo-4}
\begin{tabular}{|p{0.23\textwidth}|p{0.23\textwidth}|p{0.23\textwidth}|p{0.23\textwidth}|}
\hline
Outcome branch & Trigger in the formal search & Conditional conclusion & Next action \\\hline
supports\_mechanism & Obligations for the averaged-focusing step and the trapped-surface criterion close across the admissible metric class, and no candidate witness satisfies all assumptions & The proposed threshold-to-trapped-surface chain is internally consistent within the restricted model class; the exclusion of stable non-black-hole remnants above the threshold is supported, not proved & Replicate the closure under independently confirmed assumptions and test the most consequential declared boundary condition \\\hline
partial\_or\_heterogeneous & The chain closes only for a sub-class of metrics or only for part of the compactness range, with variation across the declared controls & The relation is conditional or heterogeneous; no universal exclusion is warranted and the validity domain must be narrowed & Predefine and check plausible moderators and improve coverage across the declared conditions and comparability of the formal constructs \\\hline
null\_or\_contradictory & A proposed witness satisfies every assumption, including the averaged-focusing condition and the stability predicate, yet evokes no trapped surface, or a proof obligation shows the assumptions cannot yield the bound & The proposed relation is not supported within this design boundary; absence of support is not proof of absence in general & Audit construct validity and comparison adequacy, then revise the mechanism or boundary claim before another design iteration \\\hline
uninformative\_or\_invalid & A required formal dependency remains undefined, so the focusing term or the compactness threshold cannot be evaluated for the candidate witnesses & No scientific conclusion is warranted because the planned analysis did not yield interpretable evidence & Resolve the identified validity failure before repeating the design and obtain human confirmation of the missing formal prerequisites \\\hline
\end{tabular}
\end{table*}

\subsection{Separating the Theorem from the Astrophysical Gap}
The matrix deliberately isolates the formal claim from its astrophysical application. Even a supportive branch constrains only the restricted model class; it does not by itself explain the observed lower mass gap, which remains sensitive to explosion energetics, fallback, equation-of-state physics, pair-instability processes, and binary evolution. A null or heterogeneous branch likewise does not refute the existence of a compactness bound in nature; it refutes the sufficiency of the declared assumptions for that bound. This separation is the section's contribution: it converts the proof search into a transparent set of conditional conclusions, so that the next stage can inherit a precise dependency rather than an ambiguous verdict.

\section{Risks, Limitations, and Human Review Requirements}
\subsection{Formal Scope and Empirical Separation}
The proposed AANEC-relaxed trapped-surface theorem operates strictly within classical general relativity under spherical, non-rotating collapse with isotropic pressure. Its validity is contingent upon global hyperbolicity, regular matter fields, complete null generators, and an explicit averaged null-convergence condition. These geometric and causal constraints define the theorem's boundary; they do not constitute a quantitative derivation of the observed 2-5 solar mass gap. The mass gap remains an astrophysical consistency constraint governed by supernova explosion energetics, fallback dynamics, equation-of-state physics, and binary evolution. Consequently, formal claims regarding trapped-surface formation must remain analytically distinct from empirical population inferences. The theorem's exclusion of stable non-black-hole remnants applies only to equilibrium objects satisfying the bounded compactness limit within this restricted model class, leaving open the possibility that nonspherical, rotating, or quantum-stress-energy-dominated configurations evade the averaged-focusing assumptions. This section owns the detailed review criteria for these limitations.

\begin{table*}[htbp]
\centering
\footnotesize
\setlength{\tabcolsep}{3pt}
\renewcommand{\arraystretch}{1.12}
\caption{Limitation and Review Decision Matrix}
\label{tab:risk-limitations-and-review-rlr-02}
\begin{tabular}{|p{0.23\textwidth}|p{0.23\textwidth}|p{0.23\textwidth}|p{0.23\textwidth}|}
\hline
Decision Node & Trigger Condition & Required Action & Release Constraint \\\hline
**A5 Formalization** & Explicit averaged null-convergence condition lacks mathematical form distinguishing it from standard AANEC. & Define the integral constraint on the stress-energy tensor along achronal null generators. & Proof obligation PO1 remains unresolved; counterexample verification is suspended until A5 is specified. \\\hline
**D4 Compactness** & Misner-Sharp compactness parameter C is undefined for the chosen metric class. & Specify the formula for effective mass within areal radius r. & Threshold C\_max cannot be calculated; exclusion bound P2 cannot be numerically instantiated. \\\hline
**A6 Stability** & Stability predicate for non-black-hole remnants is ambiguous relative to regular black hole models. & Define the TOV-like equilibrium limit that distinguishes stable remnants from horizon-bearing objects. & Counterexamples CE1/CE2 cannot be adjudicated as valid physical witnesses or definitional conflicts. \\\hline
**Formal Review** & The final canonical formal\_reasoning\_plan status requires qualified human review. & Expert verification of the Raychaudhuri focusing chain and trapped-surface identification criterion. & The theorem remains a conjecture; no claim of proof or established exclusion may be published. \\\hline
\end{tabular}
\end{table*}

\subsection{Counterexample Adjudication and Proof Obligations}
The search for counterexamples to the exclusion proposition P2 has identified two candidate witnesses, CE1 and CE2, both currently classified as unverified boundary cases. CE1 posits a static, high-compactness regular metric that may evade trapped surfaces, while CE2 invokes regular black hole models with non-singular cores. However, the adjudication of these witnesses is blocked by the missing formal definition of assumption A5. If A5 is interpreted as a strict integrated focusing constraint independent of AANEC, CE1 likely violates this assumption, rendering it an invalid counterexample to the theorem as stated. Similarly, CE2's status as a 'stable non-black-hole remnant' is compromised if the presence of a horizon in regular black hole models contradicts the A6 stability predicate. Until the mathematical form of A5 and the precise definition of C in D4 are supplied, proof obligation PO1 cannot be closed. The current analysis does not demonstrate that AANEC plus averaged focusing yields a trapped-surface bound; it only identifies the specific formal gaps that prevent falsification or verification. This section details the human-review requirements for these specific formal gaps.

\subsection{Alternative Mechanisms and Boundary Failures}
The theorem's explanatory relevance is limited by several declared alternative mechanisms. Pointwise energy conditions are systematically violated by quantum fields, necessitating the use of averaged conditions for broader validity. However, AANEC may hold while local energy-condition violations or quantum stress-energy effects prevent the pointwise focusing required for trapped-surface formation. Furthermore, the mass-gap inference fails to discriminate the theorem from astrophysical alternatives if explosion energetics, pair-instability processes, or binary evolution dominate the observed remnant distribution. The lower edge of the black hole mass gap is dictated by near-final core mass and is robust against variations in metallicity and internal mixing, suggesting that stellar physics, rather than pure GR compactness bounds, sets the observed demarcation. These alternative explanations do not falsify the formal theorem but restrict its applicability to the idealized spherical, isotropic, and classical regime defined by the assumptions. This section owns the detailed limitation decision ledger regarding these alternative mechanisms.~\cite{cite_0ed84243cf171bc4,cite_7a407c865faafd11,cite_789b47fb0241d9ae}

\subsection{Concrete Human-Review Decisions}
\paragraph{Human-review checklist.} 1. **Supply A5 Definition**: Provide the exact integral expression for the explicit averaged null-convergence condition to distinguish it from AANEC (A2) and enable verification of witness CE1. 2. **Define D4 Metric**: Specify the Misner-Sharp compactness formula for the spherically symmetric metric class to calculate C\_max and the exclusion threshold. 3. **Clarify A6 Stability**: Define the stability predicate for non-black-hole remnants to determine if regular black hole models (CE2) satisfy or violate the premise of being 'stable' and 'non-black-hole'. 4. **Review Formal Plan**: Conduct qualified human review of the formal\_reasoning\_plan status to confirm that the Raychaudhuri derivation steps S1-S4 remain proposed and unverified.

\section{Definitions, Propositions, and Proof Obligations}
\subsection{Scope of the Formal Model}
This section establishes the compactness conventions and symbol domains required for the AANEC-relaxed trapped surface argument. The model restricts attention to classical general relativity under spherically symmetric collapse with isotropic pressure, ensuring that all subsequent derivations remain within a well-defined geometric class. By isolating these structural prerequisites, we prevent the conflation of formal geometric bounds with astrophysical explanations for the observed stellar mass gap.

\subsection{Core Symbol Definitions}
\paragraph{Definition.} The following definitions provide the operational basis for the geometric and matter variables used throughout the proof obligations:

1. Mass of the collapsing matter or remnant object. 2. Areal radius or radial coordinate. 3. Expansion scalar of a null geodesic congruence. 4. Stress-energy tensor. 5. Spacetime metric. 6. Tangent vector to null geodesic generators.

\subsection{Misner-Sharp Compactness Parameter}
\paragraph{Definition.} The Misner-Sharp compactness parameter is defined as the ratio of the effective mass enclosed within an areal radius to that radius. This parameter serves as the primary geometric threshold indicator for trapped surface formation.

\begin{equation}
C = 2G M_{eff} / (c^2 r)
\label{eq:definitions-and-propositions-eq-compactness}
\end{equation}

\subsection{Trapped Surface Formation Threshold}
\paragraph{Proposition (proposed).} Proposition P1: In spherically symmetric stellar collapse satisfying AANEC, global hyperbolicity, regular matter fields, complete null generators, and an explicit averaged null-convergence/focusing condition, reaching a model-dependent compactness threshold forces trapped-surface formation.

\subsection{Exclusion of Stable Non-Black-Hole Remnants}
\paragraph{Proposition (proposed).} Proposition P2: Stable non-black-hole remnants above a compactness-dependent mass scale are excluded within the restricted model class defined by AANEC and averaged focusing.

\subsection{Unresolved Proof Obligations}
\begin{itemize}
\item The formal argument currently rests on three unresolved proof obligations that must be closed to validate the proposed propositions:
\item PO1: Establishing the precise mathematical form of the averaged null-convergence condition and verifying that candidate witnesses satisfy it.
\item PO2: Demonstrating that the proposed trapped surface identification criterion holds for all admissible metrics within the restricted model class.
\item PO3: Deriving the explicit exclusion bound for stable non-black-hole remnants from the compactness threshold and TOV-like stellar-structure assumptions.
\end{itemize}

\begin{table*}[htbp]
\centering
\footnotesize
\setlength{\tabcolsep}{3pt}
\renewcommand{\arraystretch}{1.12}
\caption{Proof Obligation Status Matrix}
\label{tab:definitions-and-propositions-table-decisions}
\begin{tabular}{|p{0.23\textwidth}|p{0.23\textwidth}|p{0.23\textwidth}|p{0.23\textwidth}|}
\hline
Obligation ID & Target Proposition & Current Status & Required Input \\\hline
PO1 & P1 & Unresolved & Mathematical form of A5 \\\hline
PO2 & P1 & Unresolved & Global hyperbolicity verification \\\hline
PO3 & P2 & Unresolved & Misner-Sharp compactness formula \\\hline
\end{tabular}
\end{table*}

\section{Forward Derivation and Counterexample Search Plan}
\subsection{Scoped Premises for the Focusing Chain}
This stage converts the scoped premises supplied by the definitions route into an obligation for a forward derivation. The argument consumes six declared assumptions: spherical collapse with isotropic pressure, AANEC along complete null generators, global hyperbolicity and regularity, completeness of the relevant generators, an explicit averaged null-convergence/focusing condition that is treated as an independent geometric hypothesis rather than a consequence of AANEC, and a TOV-like bounded-compactness model for stable non-black-hole remnants. Two candidate propositions are inherited: a threshold criterion forcing trapped-surface formation, and an exclusion criterion for stable remnants above a compactness-dependent mass scale. The compactness symbol and the mathematical form of the averaged focusing hypothesis are carried forward as unresolved dependencies owned by the definitions ledger; this section does not restate their full review records, but uses them as the variables whose closure determines whether the derivation can be certified. This approach advances the argument by focusing on the logical structure of the derivation obligations rather than repeating the definition status.

\subsection{Working Convention for the Focusing Parameter}
\paragraph{Definition.} For the derivation to be testable, the averaged focusing hypothesis must be expressed as a signed convergence functional along a null generator. We introduce a working convention: let the averaged focusing parameter be the affine-length-normalized integral of the convergence term over a complete generator segment, required to be non-negative and sufficiently large to overcome any initial positive expansion. This convention is a proposed operationalization, not a completed definition; its precise integrand and the threshold value that qualifies as "sufficiently large" are the dependencies that the definitions ledger must fix before the chain can be closed.

\subsection{Integrated Focusing Relation}
\begin{equation}
theta(1) <= theta(0) - (1/2) * theta(0)^2 - sigma
\label{eq:forward-derivation-and-counterexamples-fd-eq}
\end{equation}

\subsection{Reading the Integrated Relation}
The displayed relation is the proposed integrated form of the Raychaudhuri chain that the derivation must justify. The expansion scalar at the end of a unit-affine segment is bounded above by its initial value minus a quadratic self-focusing term and minus the averaged focusing parameter. The obligation is to show that, under the declared assumptions, a non-negative and sufficiently large averaging functional drives the expansion negative within finite affine parameter, so that the congruence converges before any regular equilibrium configuration can be maintained. Because the integrand defining the averaging functional is not yet fixed, the relation is stated as a proof target rather than an established identity. This section advances the argument by specifying the logical form of the derivation obligation, distinct from the definition status owned by the previous section. See Eq.~\eqref{eq:forward-derivation-and-counterexamples-fd-eq}.

\subsection{Lemma-by-Lemma Derivation Obligations}
\begin{itemize}
\item Lemma 1 (Focusing): from the integrated relation, establish that a complete generator satisfying the averaged focusing hypothesis reaches negative expansion within bounded affine parameter; the failure condition is an incomplete generator or an averaging functional that admits arbitrarily small positive values.
\item Lemma 2 (Trapped-surface identification): from negative expansion of the outgoing congruence together with the ingoing congruence fixed by spherical regularity, identify a closed two-surface on which both expansions are negative; the failure condition is a metric whose compactness exceeds the threshold without producing simultaneous sign reversal of both expansions.
\item Lemma 3 (Equilibrium incompatibility): show that a trapped surface is inconsistent with the bounded-compactness equilibrium model assumed for stable remnants; the failure condition is a stable configuration whose structure evades the compactness bound.
\item Lemma 4 (Exclusion bound): combine the trapped-surface threshold with the mass definition to express a mass scale above which no stable non-black-hole remnant can exist within the model class; the failure condition is an unresolved mapping between the compactness threshold and the mass variable.
\end{itemize}

\subsection{Proposed Counterexample Criterion}
\paragraph{Proposition (proposed).} A proposed criterion for an admissible counterexample: a witness falsifies the exclusion conclusion only if it simultaneously satisfies all six scoped assumptions and yet fails to form a trapped surface above the compactness threshold. A witness that violates the averaged focusing hypothesis, the completeness of generators, or the bounded-compactness equilibrium model is an out-of-domain boundary case, not a refutation. This criterion is unverified and depends on the human-supplied mathematical form of the focusing hypothesis.

\begin{table*}[htbp]
\centering
\footnotesize
\setlength{\tabcolsep}{3pt}
\renewcommand{\arraystretch}{1.12}
\caption{Counterexample Decision Matrix}
\label{tab:forward-derivation-and-counterexamples-fd-matrix}
\begin{tabular}{|p{0.23\textwidth}|p{0.23\textwidth}|p{0.23\textwidth}|p{0.23\textwidth}|}
\hline
Candidate witness & Failing assumption & Decision & Consequence for the theorem \\\hline
Static high-compactness regular metric with isotropic pressure and AANEC but evading trapped surfaces & Averaged focusing hypothesis and bounded-compactness equilibrium & Out-of-domain boundary case & Does not refute the theorem; shows the added focusing assumption is doing the work that AANEC alone cannot \\\hline
Regular black-hole geometry with a resolved core but a horizon & Stable non-black-hole remnant predicate (the object is a black hole, not a horizonless remnant) & Assumptions not satisfied & Confirms the trapped-surface step rather than challenging the exclusion of horizonless remnants \\\hline
Object satisfying every scoped assumption including the fixed focusing functional, yet no trapped surface above threshold & None & Valid counterexample & Forces restriction or rejection of the exclusion conclusion within the model class \\\hline
Object whose generators are incomplete or whose matter violates AANEC & Completeness or AANEC & Out-of-domain boundary case & Narrows the validity domain without altering the theorem inside its declared scope \\\hline
\end{tabular}
\end{table*}

\subsection{Limitation of the Derivation Plan}
The plan cannot presently be certified because two of its load-bearing symbols are unresolved: the compactness parameter and the integrand of the averaged focusing hypothesis. Until these are fixed, every candidate witness is judged against a moving definition of the assumptions, and the decision matrix can only classify boundary cases rather than produce a definitive counterexample. A further limitation is the deliberate separation of the formal claim from the astrophysical mass-gap inference: even a closed derivation would not by itself explain the observed distribution of compact-object masses, which remains subject to explosion energetics, fallback, equation-of-state physics, rotation, magnetic fields, nonspherical dynamics, and binary evolution. The mass gap is therefore treated as a consistency constraint that the theorem must not contradict, not as a prediction that the theorem supplies. This section hands the unresolved dependencies and the decision criteria to the outcomes stage, where the prespecified analysis will register whether the derivation obligations close, remain conditional, or are invalidated by an admissible witness. This approach advances the argument by focusing on the derivation's logical dependencies rather than repeating the definition ledger's status.~\cite{cite_0ed84243cf171bc4,cite_7a407c865faafd11}

\appendices
\section{Idea Source Checkpoints and Direction Selection Audit}
\subsection{Audit Snapshot Inventory and Selection Priority}
The selected direction, mechanism replacement, is reconstructed from three available audit snapshots rather than a temporal iteration series. Because no timestamps, parent identifiers, or auditable diffs are supplied, the snapshots are treated as coexisting audit records with unknown chronology. The selection priority places the selected direction first, followed by the legacy best entry, the selected primary idea, and the candidate idea. This ordering determines which formulation is carried forward: the retained decision is the AANEC-relaxed trapped-surface bound, while the observed stellar mass gap is demoted from a theorem-derived prediction to an astrophysical consistency constraint. The snapshots do not constitute empirical evidence for the formal claim; they record how the intellectual object was framed, scoped, and prioritized upstream.

\begin{table*}[htbp]
\centering
\footnotesize
\setlength{\tabcolsep}{3pt}
\renewcommand{\arraystretch}{1.12}
\caption{Retained Decisions and Scope Corrections}
\label{tab:appendix-idea-evolution-block-2}
\begin{tabular}{|p{0.29\textwidth}|p{0.29\textwidth}|p{0.29\textwidth}|}
\hline
Audit checkpoint & Retained decision & Scope correction carried into the route \\\hline
Candidate snapshot & Replace pointwise NEC/SEC with AANEC plus an explicit averaged null-convergence/focusing condition. & The added focusing condition is declared independent of AANEC, so AANEC alone is not treated as sufficient for trapped-surface formation. \\\hline
Selected primary idea & Keep the bound conditional on a model-dependent compactness threshold. & The bound is not a quantitative derivation of the 2-5 solar mass gap and is not a general proof that all massive stars form black holes. \\\hline
Direction hypothesis & Demote the mass gap to an astrophysical consistency constraint. & Explosion energetics, fallback, equation-of-state physics, rotation, magnetic fields, nonspherical dynamics, and binary evolution remain admissible alternative explanations that can weaken explanatory relevance without falsifying the formal theorem. \\\hline
\end{tabular}
\end{table*}

\subsection{Open Questions Preserved for Review}
Three unresolved items are preserved rather than silently closed, because they constrain how the retained decision can be operationalized. First, the exact mathematical form of the explicit averaged null-convergence condition that distinguishes it from standard AANEC is not fixed by any snapshot, so the added assumption cannot yet be checked against a proposed witness. Second, the Misner-Sharp compactness parameter lacks a formal definition for the chosen metric class, leaving the compactness threshold and the exclusion bound underdetermined. Third, the stability predicate for non-black-hole remnants is not separated from regular black hole constructions, so a candidate object may satisfy the theorem's premises only to be reclassified as a black hole rather than an excluded remnant. These gaps are recorded as open questions for human review; they are not treated as completed proofs, observed findings, or settled definitions. This section advances the proposal by documenting the provenance of these dependencies, while the detailed resolution criteria remain owned by the definitions and review ledgers.

\subsection{Transition to the Formal Section}
The audit therefore hands off a bounded proposal: a conditional trapped-surface/compactness theorem under AANEC plus an explicit averaged-focusing assumption, with the mass gap retained only as a consistency check. The formal section should take the retained decision as its premise and the three open questions as its dependency list. It must state the averaged-focusing condition and the compactness parameter as proposed formal objects, record the proof obligations as unverified, and treat any candidate counterexample as invalid unless it satisfies every declared assumption, including the independent focusing condition and the stability predicate. No outcome branch is observed; the design unit remains a formal proof and counterexample search.

\section{Variables, Symbols, and Operational Definitions}
\subsection{Core Geometric and Matter Symbols}
The formal model class is parameterized by a set of geometric and matter variables. The spacetime metric g and stress-energy tensor T define the background geometry and matter content. Null geodesic generators are characterized by the tangent vector k and the expansion scalar theta. The AANEC is operationally defined as the integral of T\_ab k\textasciicircum{}a k\textasciicircum{}b along null geodesics k\textasciicircum{}a being non-negative. The Misner-Sharp compactness parameter C and the mass of the collapsing matter or remnant object M are central to the threshold criteria. These symbols form the basis for the proposed Raychaudhuri-type focusing and trapped-surface identification.~\cite{cite_0ed84243cf171bc4}

\begin{table*}[htbp]
\centering
\footnotesize
\setlength{\tabcolsep}{3pt}
\renewcommand{\arraystretch}{1.12}
\caption{Unresolved Operational Definitions}
\label{tab:appendix-variables-and-definitions-block-2}
\begin{tabular}{|p{0.23\textwidth}|p{0.23\textwidth}|p{0.23\textwidth}|p{0.23\textwidth}|}
\hline
Variable & Operational Role & Status & Decision Required \\\hline
V6 & Misner-Sharp compactness C & needs\_human\_input & Define formula for effective mass within radius r in chosen metric. \\\hline
V7 & Stellar remnant mass M & needs\_human\_input & Specify measurement method (e.g., radial velocity curves vs pulsar timing). \\\hline
V8 & Explosion energy & needs\_human\_input & Define specific metric (kinetic vs total) and measurement in observed samples. \\\hline
V10 & Equation of state models & needs\_human\_input & List specific models or nuclear physics parameters to be tested. \\\hline
V11 & Rotation & needs\_human\_input & Define measurement or initial condition specification for collapse models. \\\hline
V12 & Magnetic fields & needs\_human\_input & Define measurement or initial condition specification for collapse models. \\\hline
\end{tabular}
\end{table*}

\subsection{Focusing Condition and Threshold Dependencies}
The explicit averaged null-convergence/focusing condition (A5) is an independent geometric hypothesis, not a consequence of AANEC alone. Its exact mathematical form distinguishing it from standard AANEC (A2) is currently missing and requires human input to resolve proof obligation PO1. Similarly, the specific formula for the Misner-Sharp compactness threshold dependent on model parameters (V6) is not defined in the brief. These unresolved operationalizations directly impact the proposed derivation of the exclusion bound (S4) and the consistency check against the stable equilibrium limit (S3). This section advances the proposal by mapping these symbolic dependencies to their downstream derivation obligations, without restating the detailed definition ledger status.

\subsection{Transition to Forward Derivation}
The completion of these operational definitions is a prerequisite for advancing the forward derivation chain. Specifically, the proposed application of the Raychaudhuri focusing equation (S1) and the trapped surface identification criterion (S2) cannot be rigorously evaluated until the mathematical form of A5 and the precise definition of C are established. The next section will develop the Raychaudhuri/integrated-threshold chain contingent upon resolving these symbolic dependencies. This section advances the proposal by explicitly linking the appendix's symbolic inventory to the forward derivation's logical requirements, ensuring that the dependency consequence is clear without repeating the definition ledger's full content.

\section{Evidence Coverage, Unknown Items, and Review Checklist}
\begin{table*}[htbp]
\centering
\footnotesize
\setlength{\tabcolsep}{3pt}
\renewcommand{\arraystretch}{1.12}
\caption{Evidence Coverage and Bounded Use}
\label{tab:appendix-evidence-and-review-block-1}
\begin{tabular}{|p{0.23\textwidth}|p{0.23\textwidth}|p{0.23\textwidth}|p{0.23\textwidth}|}
\hline
Claim Category & Source Evidence & Scope Boundary & Release Criterion \\\hline
AANEC Necessity & W3009127927 & Theoretical necessity for averaged conditions due to quantum violations of pointwise conditions. & Formal claim must remain separated from empirical observation. \\\hline
Mass Gap Origin & W2994776863 & Astrophysical mechanism (pair-instability) explains observed gap; not a direct GR theorem consequence. & Theorem must be explicitly qualified to avoid conflating GR bounds with stellar physics. \\\hline
\end{tabular}
\end{table*}~\cite{cite_0ed84243cf171bc4,cite_7a407c865faafd11}

\subsection{Unknown-Item Consequences}
\begin{itemize}
\item The current formal proposal cannot be advanced to proof status because two critical mathematical definitions remain undefined. The absence of the explicit formula for the Misner-Sharp compactness parameter (D4) prevents the calculation of any numerical threshold for the theorem's conclusion. Furthermore, the mathematical form of the explicit averaged null-convergence condition (A5) is missing, which makes it impossible to verify whether the proposed counterexamples genuinely violate the theorem's assumptions or merely reflect ill-defined premises. These unknowns restrict the current work to a conceptual framework rather than a rigorous derivation. This section advances the proposal by documenting the evidentiary status of these unknowns, while the detailed review criteria remain owned by the risk and limitations section.
\end{itemize}

\subsection{Release Criteria for Future Claims}
Future publication or release of this theoretical work requires the resolution of the formal reasoning plan's status. Specifically, the proposed theorem and its associated counterexamples must undergo qualified human review to ensure that the added averaged-focusing assumption is mathematically distinct from standard AANEC and that the stability predicates for non-black-hole remnants are rigorously defined. Until these formal specifications are provided and reviewed, all claims regarding the exclusion of stable remnants must remain marked as unverified proposals. This section advances the proposal by stating the release criteria for the appendix's evidence, while the detailed human-review decision ledger remains owned by the risk and limitations section.\section*{References}
\bibliographystyle{IEEEtran}
\bibliography{references}
\end{document}
